\subsection{点的量子 K理论}%
\label{sub:Quantum K-theory of a point}

一个自然的问题是研究高亏格量子 K理论。我们从最简单的目标，即一个点，开始。研究高亏格的一个动机是尝试证明 Calabi-Yau 四维簇的 Klemm-Pandharipande 定义的 GV 不变量或 Calabi-Yau 五维簇的 Pandharipande-Zinger 定义的 GV 不变量的整性（在这种情形下，我们只需考虑亏格 \(1\)）。我们需要考虑点的量子 K理论的一个原因在于，在 \(I\bar{M}_{1,1}(X, \beta)\) 中，存在一些连通分支，其中整条亏格 \(1\) 的曲线被收缩为一点。

\begin{quest}
    K理论版本的 Witten-Kontsevich 定理是什么？\footnote{据 You-Cheng 说，困难无处不在。}
\end{quest}

我们先回顾上同调版本的 Witten 猜想。我们从算子
\[ L = \partial_x^2 + u(x) \]
出发，并将其形变为
\[ \tilde{L} = \partial_x^2 + u(x, T_3, T_5, \ldots), \]
其中我们应将 \(x = T_1\) 且 \(u(x) = u(x, 0, 0,\ldots)\)。形式地，利用性质
\[ \partial^n \circ f = \sum_{m = 0}^{\infty} \binom{n}{m} f^{(m)} \partial^{n-m} \]
（对所有整数 \(n \in \Z\) 成立），将 \(L^{1/2}\) 写成
\[ L^{1/2} = \partial_x + \sum_{k > 0} a_{-k}(x) \partial_x^{-k}. \]
特别地，我们由方程
\[ \pdv{\tilde{L}}{T_{2k+1}} = \ab[(\tilde{L}^{\frac{2k+1}{2}})_{\geq 0}, \tilde{L}] \]
（对所有 \(k \geq 0\)）构造（2-）KdV 族。我们把满足
\[ 2 \pdv[2]{}{x}\log \tau(x, T_3, T_5, \ldots) = u(x, T_3, T_5, \ldots) \]
的函数 \(\tau\) 称为 \textit{\(\tau\)-函数}。现在用 \(u = 2x\) 进行上述过程。定义函数
\[ F(t_0, t_1, \ldots) \coloneqq \sum_{n \geq 0} \sum_{k_0, \ldots, k_n} \ab<\tau_0^{k_0} \cdots \tau_n^{k_n}> \frac{t_0^{k_0} \cdots t_n^{k_n}}{k_0! \cdots k_n!}. \]
注意到由于亏格完全由插入的个数和插入的次数决定，记号
\[ \ab<\tau_0^{k_0} \cdots \tau_n^{k_n}> \]
唯一确定了一个点的后代 GW 不变量。

\begin{thm}[Kontsevich {\cite{Kontsevich1992}}]
    \(F\) 是从 \(L = \partial_x^2 + 2x\) 出发构造的 KdV 族的 \(\tau\)-函数，其中变量替换为 \(t_i = (2i+1)!! T_{2i+1}\)。
\end{thm}

这个定理的一个应用是考虑虚局部化。例如，假设我们考虑 \(\bar{M}_{g,0}(X, d[C])\)，其中 \(C\) 是 Calabi-Yau 三维簇中唯一的有理曲线，法丛为 \(N = \mc{O}(-1)\oplus \mc{O}(-1)\)。那么在此设定下我们可以化归到局部模型，并通过虚局部化计算
\begin{align*}
    \int_{[\bar{M}_{g,0}(\P^1, d)]^{\vir}} c_{\mr{top}}(R^1 \pi_* N) = \sum_{m \vdash d} \frac{(-1)^{d-L(m)}}{\abs{\Aut(m)}} \int_{\bar{M}_{g, L(m)}} \frac{1 + \cdots + c_g(\mc{H})}{\prod_i (1-m_i \psi_i)}
\end{align*}
。如果我们希望在 K理论中得到类似的结果，那么就需要考虑形如
\[ \chi\ab(\bar{M}_{g, L(m)}/\Aut(m), \frac{\lambda_{-1}(\mc{H})}{1-m_i L_i}) \]
的项，因此我们需要考虑点的置换等变量子 K理论。

下面我们进行文献综述。首先，Lee~\cite{Lee1997} 计算了公式
\[ \ab<\frac{1}{1-q_1 L}, \ldots, \frac{1}{1-q_n L}>_{0,n}^{\mr{pt}, K} = \ab(\prod_{i=1}^n \frac{1}{1-q_i}) \ab(1 + \sum_{i=1}^n \frac{q_i}{1-q_i})^{n-3}. \]
后来，Lee--Qu~\cite{LeeQu2014} 给出了一种原则上计算
\[ \ab<\frac{1}{1-q_1 L}, \ldots, \frac{1}{1-q_n L}>_{1,n}^{\mr{pt}, K} \]
的方法。Givental~\cite{Givental2017K1} 计算了公式
\[ \sum_{n \geq 2} \ab<\frac{1}{1-q L}, 1, \ldots, 1>_{0,1+n}^{S_n} = (1-q) \exp\ab(\sum_{k > 0} \frac{\Psi^k}{k} \ab(1 + \frac{q}{1-q})). \]
最后，Milanov~\cite{Milanov2025} 在一份 2024 年首次公开的预印本中计算了
\begin{align*}
    \sum_{n \geq 0}& \ab<\frac{1}{1-q_1 L}, \ldots, \frac{1}{1-q_r L}, 1, \ldots, 1>_{0,r+n}^{S_n} \\
    &= \ab(\prod_{i=1}^{r} \frac{1}{1-q_i})\ab(1 + \sum_{i=1}^r \frac{q_i}{1-q_i}) \exp\ab(\sum_{k > 0} \frac{\Psi^k}{k} \ab(1 + \sum_{i=1}^r \frac{q_i}{1-q_i})). 
\end{align*}

以下多重霍奇递推以及基于它们的新的亏格一计算出自 You-Cheng Chou，并在本讲中给出。

设 \(s \in \Z\)，定义多重霍奇层 \(\mc{H}_s \coloneqq R^0 \pi_* (\omega_{\pi}^s)\)。考虑映射 \(\bar{M}_{g,n+1} \xrightarrow{\pi} \bar{M}_{g,n}\)，我们将上方的余切线记为 \(L_i\)，下方的记为 \(\ell_i\)。

\begin{lem}
    对所有整数 \(s\) 和 \(d_i\)，我们有
    \[ \pi_*\ab(\omega_{\pi}^s\ab(\sum_{i=1}^n d_i D_i)) = \mc{H}_s-\mc{H}_{1-s}^* + \sum_{i=1}^n \sum_{k=1}^{d_i} \ell_i^{-k+s}. \]
\end{lem}

\begin{rmk}
    注意到若 \(f(q) \in \Z[q]\)，则 \(\exp \sum_k \frac{\Psi^k}{k} (f(q)) \in \Z \ps{q}\)。原因在于这一运算把 \(f = \sum_{i=1}^{n} a_i q^i\) 映为 \(\prod_{i=1}^n \frac{1}{(1-q^i)^{a_i}}\)。
\end{rmk}

\begin{proof}
    当所有 \(d_i\) 都为零时，由 Serre 对偶我们看到
    \[ R^0 \pi_* \omega_{\pi}^s - R^1 \pi_* \omega_{\pi}^s = \mc{H}_s - \mc{H}_{1-s}^* \]
    。下一步是证明
    \[ \pi_* \omega_{\pi}^s ((r+1)D_i) = \pi_* \omega_{\pi}^s(r D_i) + \ell_i^{s-r-1}. \]
    这里我们先考虑正合列
    \[ 0 \to \mc{O} (-D_i) \to \mc{O} \to \mc{O}_{D_i} \to 0 \]
    并用 \(\omega_{\pi}^s\) 扭变得到
    \[ 0 \to \omega_{\pi}^s(r D_i) \to \omega_{\pi}^s ((r+1)D_i) \to \omega_{\pi}^s ((r+1)D_i) \otimes \mc{O}_{D_i} \to 0 \]
    ，然后计算前推。
\end{proof}

\begin{lem}
    设 \(f_i\) 为有理函数，例如 \(\frac{1}{1-q L}\)。则
    \begin{align*}
        \ab<f_1(L), \ldots, f_n(L), L_{n+1}^r>_{g, n+1} ={}& \sum_{i=1}^n \ab<f_1(L), \ldots, D(f_i f_{n+1}), \ldots, f_n(L)>_{g, n} \\
        &+ \ab<f_1(L), \ldots, f_n(L) \mid \mc{H}_r - \mc{H}_{1-r}^*>_{g,n},
    \end{align*}
    其中 \(f_{n+1} = L_{n+1}^r\)，且 \((Df)(x) \coloneqq \frac{f(x)-f(1)}{x-1}\)。
\end{lem}

为证明这一点，我们需要考虑前推
\begin{align*}
    \pi_* (L_1^{d_1} \cdots L_{n+1}^{d_{n+1}}) &= \pi_* \ab(\pi^*(\ell_1^{d_1} \otimes \cdots \otimes \ell_n^{d_n}) \otimes \omega_{\pi}^{d_{n+1}}\ab(\sum_i (d_i+d_{n+1})D_i)).
\end{align*}
这里我们用到 \(L_i = (\pi^* \ell_i)(D_i)\) 与 \(L_{n+1} = \omega_{\pi}(D_1 + \cdots + D_n)\) 这两个事实。应用上一个引理，我们看到前推变为
\[ (\ell_1^{d_1} \cdots \ell_n^{d_n})\ab(\mc{H}_{d_{n+1}} - \mc{H}_{1-d_{n+1}}^* + \sum_{i=1}^n \sum_{k=1}^{d_i + d_{n+1}} \ell_i^{-k+d_{n+1}}). \]
现在，Lee 的公式通过展开式
\[ \frac{1}{1-qL} = \frac{1}{1-q} \sum_{k \geq 0} \frac{q^k(L-1)^k}{(1-q)^k} \]
计算关联子
\[ \ab<(L-1)^{k_1}, \ldots, (L-1)^{k_n}>_{0,n} = \frac{(n-3)!}{\prod_{i=1}^n k_i! \ab(n-3-\sum_i k_i)!} \]
。这里注意在亏格 \(0\) 没有霍奇丛，因此实际上可以用该引理导出 Lee 的公式。

\begin{lem}
    我们有
    \[ \mc{H}_r - \mc{H}_{1-r}^* = \pi^* (\mc{H}_r - \mc{H}_{1-r}^*) - \sum_{i=1}^n (D^2 (\pi^* \ell_i)^r) \mc{O}_{D_i}. \]
\end{lem}

这些引理合起来表明，我们可以化归到只有一个标记的情形，代价是插入项变得非常复杂，并且处处引入多重霍奇项。在亏格 \(1\)，事情变得容易，因为在 \(\bar{M}_{1,1}\) 上我们有
\[ \mc{H}_r = \mc{O}(r) = L_1^r. \]
此时，回顾经典公式
\[ \chi\ab(\bar{M}_{1,1}, \frac{1}{1-qL}) = \frac{1}{(1-q^4)(1-q^6)}. \]
遗憾的是，计算
\[ \chi\ab(\bar{M}_{1,n}; \frac{1}{1-q_1 L} \cdots \frac{1}{1-q_n L}) \]
太困难了，因此我们可以尝试计算
\[ \chi\ab(\bar{M}_{1,1}; \frac{1}{1-q_1 L_1} \cdots \frac{1}{1-q_n L_1})  \]
，这是已知的，但目前还不清楚如何为它写出一个好的公式。考虑它的原因在于运算
\[ \Z\ps{q} \xrightarrow{[]_r} \Z\ps{q_1, \ldots, q_r} \qquad q^k \to h_k(q_1, \ldots, q_r) \coloneqq \sum_{\ell_1 + \cdots + \ell_r = k} q_1^{\ell_1} \cdots q_r^{\ell_r}. \]
遗憾的是，You-Cheng 不知道如何以好的形式写出
\[ \chi\ab(\bar{M}_{1,1}; \frac{1}{1-q_1 L_1} \cdots \frac{1}{1-q_n L_1}) = \ab[\frac{1}{(1-q^4)(1-q^6)}]_n \]
。原则上我们可以写出形如
\[ [f(q)]_2 \coloneqq \frac{q_1 f(q_1) - q_2 f(q_2)}{q_1 - q_2} \]
的公式，但这似乎帮助不大。

\subsection{点的置换等变理论}%
\label{sub:Permutation-equivariant theory of a point}

我们可以用另一种方式计算 \(\chi\ab(\bar{M}_{1,1}, \frac{1}{1-q L})\)。考虑椭圆曲线上的对合，它给出一个映射
\[ \bar{M}_{1,1} \to \bar{M}_{0,1+3}/S_3. \]
于是，Givental~\cite{Givental2017K1} 计算出
\[ \chi\ab([\bar{M}_{0,1+n-1}/S_{n-1}], \frac{1}{1-qL}) = \prod_{i=2}^{n-1} \frac{1}{1-q^i}. \]
特别地，我们得到
\[ \frac{1}{(1-q^2)(1-q^3)}, \]
又因为该映射是一个 \(\mu_2\)-gerbe（或在对象层面是 \(2:1\) 覆叠），我们就得到了所要的公式。

考虑 \(\mc{H}_{g,1} \subset \bar{M}_{g,1}\) 为超椭圆轨迹。注意标记落在一个 Weierstrass 点上。注意超椭圆对合的不动点被映到 \(\P^1\) 上的 \(2g+1\) 个点，因此人们会希望用计算
\[ \chi\ab(\bar{M}_{0,1+2g+1}/S_{2g+1}, \frac{1}{1-q L}) = \prod_{i=2}^{2g+1} \frac{1}{1-q^i} \]
来计算量
\(\chi\ab(\mc{H}_{g,1}, \frac{1}{1-qL})\)。然而，由于上方曲线的退化方式，我们需要考虑最后 \(2g+1\) 个标记权重为 \(\frac{1}{2}\) 的 Hassett 模空间~\cite{Hassett2003}。
 






